\begin{houseviewbox}[投委会结论：开放研究宇宙；当前没有获批组合]
经授权修订后，本案只保留滚动五年、成本后年化TWRR达到5\%这一项条件性、事后研究目标。市场策略不能被诚实地包装为年度最低收益或收益保证；目标也不等同于概率判断或实施授权。已结算现金分配只作为总回报核算和产品尽调中的事实输入：它在总回报中只处理一次，其后留存、取出或日后另行考虑配置，均不创造第二层回报。本框架采用比例口径、与资金规模无关，不预设任何账户层事实。

原有的 ETF／REIT 产品菜单和10\%／15\%／40\%／25\%／10\%向量，只能代表首轮研究设计，不能成为永久投资禁区或隐性配置指令。它会在研究开始前排除可能具有更好风险调整回报、现金质量、估值或分散化贡献的直接股票及其他合规工具。自本版起，研究宇宙对所有合规可及的直接股票、ETF、公募REIT、债券和现金工具开放；现有十只产品仅是已完成首轮逐项尽调的队列。

截至本报告日，T3批准可投池和T4批准组合均为空：没有指定直接股票，十只产品也均未完成实时估值、流动性、账户、组合回放和风险预算放行。因此当前动作仍是\textbf{不配置任何市场交易资产}。这表示模型未授予市场风险配置权限；合格受保障流动性仅是未通过资格时的条件性回退类别，不满足收益目标。
\end{houseviewbox}

\begin{exhibitbox}[研究宇宙、可投池与组合不是同一个集合]
\small
\begin{tabularx}{\exhibitboxwidth}{L{2.2cm}L{3.0cm}P P}
\toprule
\textbf{层级} & \textbf{定义} & \textbf{当前状态} & \textbf{能否配置} \\
\midrule
T0 研究边界 & 法律、账户、税务、结算、流动性与授权允许研究的产品格式；不是代码白名单。 & 已定义。杠杆、衍生品、私募或显著低流动性工具需另行书面授权。 & 否。\\
T1 开放研究宇宙 & 可提名直接股票、ETF、公募REIT、主权债与现金工具。现有十只不是排他菜单。 & 开放；所有新提名默认$q=0$。 & 否。\\
T2 决策宇宙 & 针对某一次评审冻结的有限候选，含代码、研究问题、责任人、时间戳和证据路径。 & 十只产品构成首轮完成队列；未冻结直接股票名单。 & 否。\\
T3 批准可投池 & 已完成单项R0--R4、独立证据/估值/风险/运营复核的具体标的。 & 空。 & 仅可进入组合评审。\\
T4 批准组合 & 从T3按边际总回报、估值、风险、重叠、容量和流动性动态选择。 & 空；$q_{portfolio}=0$。 & 只有取得组合审批后才可讨论。\\
\bottomrule
\end{tabularx}
\end{exhibitbox}

\section{目标、证据与当前动作}

\begin{exhibitbox}[目标不是配额，研究通过也不是配置授权]
\small
\begin{tabularx}{\exhibitboxwidth}{L{3.0cm}P P}
\toprule
\textbf{约束／目标} & \textbf{当前证据结论} & \textbf{当前动作} \\
\midrule
每年最低5\%收益 & 未被证明；市场策略不能作最低收益承诺。研究口径为滚动五年、成本后TWRR目标。 & 不把“目标”写成承诺；正式回放完成前不批准实施。 \\
现金分配与税费处理 & 历史分派只能说明产品机制或历史事实，不能替代未来总回报、估值或实施资格。 & 将分配、税费、成本和返还性成分用于总回报归因与产品尽调；不设独立收款门槛。 \\
标的选择自由度 & 直接股票原被预先排除，限制了跨资产比较；开放宇宙只解决提名权，不代表自动投资资格。 & 建立有限、带日期的决策宇宙；每项按其结构完成尽调。 \\
实时实施资格 & 市场、账户、容量、税费、相关性与适当性资料均未完整。 & 每个候选$q=0$；没有下单、份额或建仓时间表。 \\
\bottomrule
\end{tabularx}
\end{exhibitbox}

\section{直接股票不是豁免，也不是永久禁区}

直接股票可以贡献盈利增长、股息增长、估值修复或与既有基金不同的经济暴露；但“高股息”本身不能证明现金可持续，更不能证明其优于 ETF、REIT 或现金工具。每一个拟纳入T3的直接股票必须先在T2名单中按代码登记，并拥有独立、可更新的研究档案：公司及行业证据、审计财务、驱动因素至收入/利润/EPS桥、治理和资本配置、当前价格下的熊/基/牛估值、股息/自由现金流覆盖、催化与证伪、A股自由流通量/停复牌/涨跌停/容量，以及同全部现有候选的发行人集团、行业、因子和共同压力映射。T3的已命名股票数必须与估值及券商材料逐一勾稽；没有匹配材料的股票不得进入T3。

同样的严谨度适用于非股票产品，但模板不能机械照搬。ETF需要基金规则、持仓穿透、跟踪和参考值机制，以及重大证券借贷、衍生品和执行能力；公募REIT需要项目、租户或承购方、运营、债务、资本开支、资产估值和分派桥。产品格式不能替代逐项尽调。

\section{从提名到组合的六道串联门槛}

\begin{enumerate}
  \item \textbf{授权与可及性：}确认产品、账户、税务、结算和流动性属于允许研究范围；
  \item \textbf{研究准入：}冻结有限候选并写明研究问题、数据截止日、责任人和来源计划；
  \item \textbf{单项权益研究级尽调：}完成适配产品结构的R0--R4证据、模型、审稿和渲染记录；
  \item \textbf{独立复核：}证据、估值、风险、流动性、税务和运营控制分别复核；
  \item \textbf{组合审批：}评估边际总回报、估值、分配与现金流质量、相关性、集中度、共同风险、压力损失、容量和适当性；没有强制类别配额；
  \item \textbf{持续监控：}资料过期、重大事件、估值/风险假设失效或流动性恶化，标的立即降为受限或不可增持。
\end{enumerate}

任一门槛缺失，候选保持$q=0$，模型不授予风险配置权限。快速事件研究可以缩短信息收集时间，但不能豁免估值可复算、风险聚合和独立复核。

\section{目前为什么没有批准组合}

本案当前没有一份冻结后的跨资产组合回放，也没有可比较的同日现金、成本、税费、流动性、集中度和风险预算数据。首轮十只产品的逐项研究已说明产品机制和证据缺口，但不等同于T3批准；直接股票尚未命名，更不能以“开放”一词直接进入组合。下一次策略评审必须从T1中冻结T2候选，再以同一套收益、风险和执行标准与现有产品比较。
